\section{Metodología de implementación}

\subsection{Enfoque metodológico del proyecto}
\begin{indentedsubsection}
La implementación del sistema de gestión en la nube para Perú Flores Travel Agency Tours Operador E.I.R.L. se rige bajo un enfoque de transformación digital ágil y 100\% práctico. Al tratarse de una solución basada exclusivamente en las funcionalidades estándar (Out-of-the-box) de Odoo Online (SaaS), el proyecto no requiere desarrollos de código a medida, lo que asegura una adopción inmediata, estable y libre de riesgos técnicos.

El enfoque principal es la transición operativa del modelo físico y manual ---basado en cuadernos para el registro de pasajeros, sumas repetitivas en Excel y confirmaciones por WhatsApp--- hacia un ecosistema unificado. El proyecto prioriza que la gerencia perciba el valor de la herramienta desde las primeras sesiones, automatizando inmediatamente sus cuellos de botella más críticos: el cálculo de cotizaciones y la generación de propuestas en PDF.
\end{indentedsubsection}

\subsection{Metodología ágil y enfoque ``Aprender Haciendo''}
\begin{indentedsubsection}
Dada la necesidad de que la gerencia adquiera autonomía operativa sobre la plataforma, se ha descartado un despliegue teórico tradicional. En su lugar, se adopta una metodología ágil combinada con el enfoque práctico de ``Aprender Haciendo''.

Este método facilita que el sistema se adapte de forma precisa a los flujos reales de la empresa, mediante ciclos cortos de configuración y retroalimentación inmediata. Sus características principales son:

\begin{itemize}
    \item \textbf{Entregas parciales y validación en tiempo real:} se configuran y presentan módulos específicos (como el cotizador) para que la gerencia los valide utilizando datos reales de paquetes turísticos (ej. paquetes a Cusco, Arequipa o Puno).
    \item \textbf{Capacitación aplicada:} las sesiones no se basan en manuales genéricos, sino en replicar el trabajo diario de la gerencia directamente en Odoo, como ingresar un nuevo lead, simular márgenes de ganancia (ej. 15\%, 20\%, 30\%) y generar el PDF final.
    \item \textbf{Flexibilidad operativa:} el sistema se va ajustando según la retroalimentación de la gerencia, ocultando elementos innecesarios (como el desglose del IGV por ítem en el PDF) para que la herramienta final sea limpia y exacta a lo que la agencia necesita.
\end{itemize}
\end{indentedsubsection}

\subsection{Herramientas de gestión y comunicación}
\begin{indentedsubsection}
Para asegurar el cumplimiento del cronograma, mantener una comunicación fluida y documentar los avances, el proyecto se apoya en las siguientes herramientas:

\begin{itemize}
    \item \textbf{Odoo Online (Entorno de Pruebas y Producción):} uso de una base de datos de prueba inicial para que la gerencia explore sin miedo a cometer errores, paralela a la base de datos oficial donde se realizará la carga definitiva.
    \item \textbf{Canales de comunicación directa:} uso de un grupo de WhatsApp oficial del proyecto (``Proyecto de digitalización'') para coordinaciones rápidas, envío de enlaces y confirmaciones.
    \item \textbf{Sesiones Híbridas (Presenciales y Virtuales):} ejecución de videollamadas para revisiones de avance rápidas y jornadas presenciales intensivas para la capacitación técnica en las oficinas de la agencia.
\end{itemize}
\end{indentedsubsection}

\subsection{Sprints de configuración (Fases Ágiles)}
\begin{indentedsubsection}
A diferencia de un desarrollo de software tradicional desde cero, el trabajo se organiza en iteraciones o sprints de configuración. En lugar de esperar meses para ver resultados, la gerencia participa activamente en el diseño de su propio flujo:

\begin{itemize}
    \item Primero se valida la configuración del CRM (registro de contactos y embudo de ventas).
    \item Luego se parametriza el Cotizador de Ventas, configurando los tarifarios base de hoteles y servicios para automatizar las sumas y porcentajes de ganancia.
    \item Posteriormente se habilitan los flujos de Pagos fraccionados (adelantos y saldos) y el módulo de Encuestas post-servicio.
\end{itemize}

Esta organización iterativa reduce la curva de aprendizaje, asegurando que, al finalizar el despliegue, la gerencia y su equipo de apoyo no solo cuenten con un software operativo, sino que estén plenamente capacitados para cotizar, vender y medir la satisfacción de sus clientes con total independencia.
\end{indentedsubsection}
